%% paper82_andrews_curtis_de.tex
%% Die Andrews-Curtis-Vermutung in der kombinatorischen Gruppentheorie (Deutsch)
%% Autor: Michael Fuhrmann
%% Specialist Mathematics Project, Build 166
%% Datum: 2026-03-12

\documentclass[12pt,a4paper]{amsart}
\usepackage{amsmath,amssymb,amsthm,mathtools,enumitem,booktabs,array,hyperref}
\usepackage[utf8]{inputenc}
\usepackage[T1]{fontenc}
\usepackage[ngerman]{babel}
\usepackage{geometry}
\geometry{margin=2.5cm}

%% Theorem-Umgebungen (Deutsch)
\newtheorem{theorem}{Satz}[section]
\newtheorem{satz}[theorem]{Satz}
\newtheorem{lemma}[theorem]{Lemma}
\newtheorem{korollar}[theorem]{Korollar}
\newtheorem{corollary}[theorem]{Korollar}
\newtheorem{proposition}[theorem]{Proposition}
\newtheorem{vermutung}[theorem]{Vermutung}
\newtheorem{conjecture}[theorem]{Vermutung}
\theoremstyle{definition}
\newtheorem{definition}[theorem]{Definition}
\newtheorem{definition_de}[theorem]{Definition}
\newtheorem{beispiel}[theorem]{Beispiel}
\theoremstyle{remark}
\newtheorem{bemerkung}[theorem]{Bemerkung}
\newtheorem{remark}[theorem]{Bemerkung}

\hypersetup{
  colorlinks=true,
  linkcolor=blue,
  citecolor=blue,
  urlcolor=blue
}

\begin{document}

\title{Die Andrews--Curtis-Vermutung\\
in der kombinatorischen Gruppentheorie:\\
Balancierte Präsentationen, AC-Züge\\
und Verbindungen zur $4$-Mannigfaltigkeitstopologie}

\author{Michael Fuhrmann}
\address{Specialist Mathematics Project, Build 166}
\email{specialist-maths@project.local}
\date{2026-03-12}

\begin{abstract}
Die Andrews--Curtis-Vermutung (1965) behauptet, dass jede balancierte Präsentation
$\langle x_1,\ldots,x_n \mid r_1,\ldots,r_n \rangle$ der trivialen Gruppe durch
eine endliche Folge von AC-Zügen (zyklische Permutation, Inversion, Konjugation,
Multiplikation einer Relation sowie Generatorinvertierung mit entsprechender
Aktualisierung) auf die Standardpräsentation
$\langle x_1,\ldots,x_n \mid x_1,\ldots,x_n \rangle$ reduziert werden kann.

Trotz ihrer elementaren Formulierung mit nur grundlegender Gruppentheorie hat
die Vermutung über sechzig Jahre allen Beweisversuchen widerstanden und gilt als
eines der schwersten offenen Probleme der kombinatorischen Gruppentheorie.
Wir behandeln die AC-Züge im Detail, präsentieren die Hauptkandidaten für
Gegenbeispiele (Akbulut--Kirby-Folge und Miller--Schupp-Präsentationen), diskutieren
Verbindungen zur $4$-Mannigfaltigkeitstopologie über den $s$-Kobordismussatz sowie
Computersuchergebnisse.
\end{abstract}

\maketitle

\tableofcontents

%% =============================================================================
\section{Einleitung}
%% =============================================================================

Im Jahr 1965 stellten J.~J.\ Andrews und M.~L.\ Curtis~\cite{AC65} folgendes
Problem über Gruppenpräsentationen:

\begin{vermutung}[Andrews--Curtis, 1965]\label{verm:AC}
Sei $P = \langle x_1,\ldots,x_n \mid r_1,\ldots,r_n \rangle$ eine balancierte
Präsentation der trivialen Gruppe ($n$ Erzeuger, $n$ Relationen; präsentierte Gruppe
$= \{1\}$). Dann kann $P$ durch eine endliche Folge der folgenden
\emph{Andrews--Curtis-Züge} (AC-Züge) in die Standardpräsentation
$S_n = \langle x_1,\ldots,x_n \mid x_1,\ldots,x_n \rangle$ überführt werden:
\begin{enumerate}[label=(AC\arabic*)]
  \item Ersetze $r_i$ durch $r_i^{-1}$.
  \item Ersetze $r_i$ durch $r_i r_j$ (für $i \neq j$).
  \item Ersetze $r_i$ durch $w r_i w^{-1}$ für ein beliebiges $w \in F_n$ (Konjugation).
  \item Ersetze $r_i$ durch eine zyklische Permutation von $r_i$.
  \item Ersetze $x_i$ durch $x_i^{-1}$ (mit entsprechender Aktualisierung aller $r_j$).
\end{enumerate}
\end{vermutung}

\begin{bemerkung}
Die Züge (AC1)--(AC4) sind die Standard-AC-Züge; (AC5) gehört manchmal zur
„stabilisierten" Version. Alle Züge liefern wieder eine balancierte Präsentation
derselben Gruppe. Die Vermutung behauptet, dass diese Züge alle balancierten
trivialen Präsentationen in eine einzige AC-Äquivalenzklasse verbinden.
\end{bemerkung}

%% =============================================================================
\section{Balancierte Präsentationen und die triviale Gruppe}
%% =============================================================================

\begin{definition_de}
Eine \emph{Präsentation} einer Gruppe $G$ ist ein surjektiver Homomorphismus
$\varphi \colon F_n \to G$ von einer freien Gruppe, geschrieben
$G = \langle x_1,\ldots,x_n \mid r_1,\ldots,r_m \rangle$, wobei $r_i \in F_n$
Wörter sind, deren Normalabschluss $\langle\!\langle r_1,\ldots,r_m \rangle\!\rangle
= \ker\varphi$. Die Präsentation ist \emph{balanciert}, falls $n = m$.
\end{definition_de}

\begin{beispiel}[Präsentationen der trivialen Gruppe]\label{bsp:trivial}
Folgende balancierte Präsentationen stellen die triviale Gruppe dar:
\begin{enumerate}[label=(\alph*)]
  \item $\langle x \mid x \rangle$: Standardpräsentation.
  \item $\langle a, b \mid ab, a^{-1}b \rangle$: balanciert, triviale Gruppe.
  \item $\langle a, b \mid a^2 b^{-1}, a^3 b^{-2} \rangle$: die Akbulut--Kirby-Präsentation
        $\mathrm{AK}(2)$, mutmaßlicher Gegenbeispielkandidat.
\end{enumerate}
\end{beispiel}

\begin{proposition}
Die Menge der balancierten Präsentationen der trivialen Gruppe ist unter AC-Zügen abgeschlossen.
\end{proposition}

\begin{proof}
Jeder AC-Zug erhält die Anzahl der Erzeuger und Relationen (alle Züge lassen $n$
unverändert). Die präsentierte Gruppe bleibt trivial: (AC1) invertiert eine
Konsequenz, (AC2) ersetzt $r_i$ durch ein Element in $\langle\!\langle r_1,\ldots,r_n \rangle\!\rangle$,
(AC3)/(AC4) ersetzen durch konjugierte Elemente, (AC5) ist ein freier Gruppenautomorphismus.
\end{proof}

%% =============================================================================
\section{Verhältnis zu Tietze-Transformationen}
%% =============================================================================

\begin{definition_de}[Tietze-Transformationen]
Zwei Präsentationen definieren dieselbe Gruppe genau dann, wenn man eine aus der
anderen durch eine endliche Folge von \emph{Tietze-Zügen} gewinnen kann:
\begin{enumerate}[label=(T\arabic*)]
  \item Hinzufügen einer Relation, die eine Folge der vorhandenen ist.
  \item Entfernen einer Relation, die eine Folge der übrigen ist.
  \item Hinzufügen eines neuen Erzeugers $y$ mit neuer Relation $y = w(x_i)$.
  \item Entfernen eines Erzeugers $x_i$, der in einer Relation $x_i = w(x_j)$ auftritt.
\end{enumerate}
\end{definition_de}

\begin{satz}[Tietze, 1908]
Zwei endliche Präsentationen definieren isomorphe Gruppen genau dann, wenn sie
durch eine endliche Folge von Tietze-Zügen verbunden sind.
\end{satz}

\begin{bemerkung}
Die AC-Züge sind restriktiver als Tietze-Züge: Sie erlauben kein Hinzufügen oder
Entfernen von Erzeugern (T3, T4), was die Balance $n$ ändern würde. Die AC-Vermutung
fragt, ob die eingeschränkte Menge der Züge (AC1)--(AC4) dennoch ausreicht,
um alle balancierten trivialen Präsentationen zu verbinden.
\end{bemerkung}

%% =============================================================================
\section{Mögliche Gegenbeispiele}
%% =============================================================================

\subsection{Die Akbulut--Kirby-Folge}

\begin{definition_de}[Akbulut--Kirby-Präsentationen]\label{def:AK}
Für jede ganze Zahl $n \geq 2$ ist die Akbulut--Kirby-Präsentation:
\[
  \mathrm{AK}(n) = \langle a, b \mid a^n b^{-(n-1)},\; a^{n+1} b^{-n} \rangle.
\]
\end{definition_de}

\begin{proposition}
$\mathrm{AK}(n)$ präsentiert die triviale Gruppe für alle $n$.
\end{proposition}

\begin{proof}
Aus $a^n = b^{n-1}$ und $a^{n+1} = b^n$ folgt $a \cdot b^{n-1} = b^n$, also $a = b$.
Einsetzen: $b^n = b^{n-1}$, also $b = e$ und $a = e$.
\end{proof}

\begin{beispiel}[$\mathrm{AK}(2)$ im Detail]
$\mathrm{AK}(2) = \langle a,b \mid a^2 b^{-1}, a^3 b^{-2} \rangle$.
Durch Computer wurde verifiziert, dass diese Präsentation in keiner Folge von
AC-Zügen bis zur Länge von mehreren Tausend Schritten auf die Standardpräsentation
reduziert werden kann. Sie ist der bekannteste Kandidat für ein Gegenbeispiel
zur Andrews--Curtis-Vermutung.
\end{beispiel}

\begin{bemerkung}
Die Schwierigkeit bei $\mathrm{AK}(n)$ besteht darin, dass jede AC-Reduktion
scheinbar eine „Zwischenkomplexifizierung" erfordert, bei der die Gesamtlänge
der Präsentation erheblich ansteigt, bevor sie vereinfacht werden kann.
Ob solche Sequenzen beschränkter Komplexität existieren, ist unbekannt.
\end{bemerkung}

\subsection{Miller--Schupp-Präsentationen}

\begin{definition_de}[Miller--Schupp-Präsentationen]\label{def:MS}
Für ein Wort $w \in F_2$ in Erzeugern $a, b$:
\[
  \mathrm{MS}(w) = \langle a, b \mid a^2, w a w^{-1} b^{-2} \rangle.
\]
\end{definition_de}

\begin{proposition}
$\mathrm{MS}(w)$ präsentiert die triviale Gruppe genau dann, wenn $w$ in geeignetem Sinne
einer ungeraden Potenz von $b$ im Quotienten entspricht.
\end{proposition}

\begin{bemerkung}
Miller und Schupp~\cite{MS99} identifizierten eine Familie von Präsentationen, die
die triviale Gruppe darstellen, für die aber keine kurzen AC-Reduktionsfolgen bekannt sind.
\end{bemerkung}

%% =============================================================================
\section{Verbindungen zur \texorpdfstring{$4$}{4}-Mannigfaltigkeitstopologie}
%% =============================================================================

\begin{satz}[$h$-Kobordismussatz, Smale 1960]\label{satz:hkobord}
Für $n \geq 5$: Jeder einfach zusammenhängende glatte $h$-Kobordismus
$(W; M_0, M_1)$ mit $\dim W = n+1$ ist diffeomorph zu $M_0 \times [0,1]$.
\end{satz}

\begin{satz}[$s$-Kobordismussatz, 1964]\label{satz:skobord}
Für $n \geq 5$: Ein glatter Kobordismus $(W; M_0, M_1)$ ist diffeomorph zu
$M_0 \times [0,1]$ genau dann, wenn die Inklusion $M_0 \hookrightarrow W$ eine
einfache Homotopieäquivalenz ist (d.h.\ verschwindende Whitehead-Torsion).
\end{satz}

\begin{bemerkung}[Versagen in Dimension 4]
Für $n = 4$ (also $5$-dimensionale Kobordismen) schlägt der $s$-Kobordismussatz im
Allgemeinen fehl: Es gibt $s$-Kobordismen zwischen $4$-Mannigfaltigkeiten, die keine
Produkte sind. Die Andrews--Curtis-Vermutung ist mit der Frage verbunden, ob bestimmte
$4$-dimensionale $h$-Kobordismen trivial sind.
\end{bemerkung}

\begin{satz}[AC und $4$-Mannigfaltigkeitstopologie]\label{satz:AC4mfgk}
Eine balancierte Präsentation $P = \langle x_1,\ldots,x_n \mid r_1,\ldots,r_n \rangle$
der trivialen Gruppe entspricht einem $2$-Komplex $K_P$ (dem Standardkomplex der
Präsentation). Dieser ist homotopieäquivalent zu einem Keil von $2$-Sphären
$\bigvee S^2$. Falls $P$ AC-reduzierbar ist, ist $K_P$ formal homotopieäquivalent
zu einem Punkt.
\end{satz}

\begin{definition_de}[Geometrische AC-Vermutung]
Ein $2$-Komplex $K$, der homotopieäquivalent zu einem Punkt ist (d.h.\ zusammenziehbar
ist), heißt \emph{geometrisch AC-trivial}, falls er durch eine Folge elementarer
Kollapse auf einen Punkt reduziert werden kann.
\end{definition_de}

\begin{vermutung}[Geometrische AC-Vermutung]\label{verm:geoAC}
Jeder endliche zusammenziehbare $2$-Komplex ist geometrisch AC-trivial.
\end{vermutung}

\begin{bemerkung}
Die geometrische AC-Vermutung ist schwächer als die ursprüngliche AC-Vermutung
(da letztere nicht-zusammenziehbare Zwischenräume erlaubt). Selbst diese schwächere
Version ist offen.
\end{bemerkung}

%% =============================================================================
\section{Der Boone--Novikov-Satz und Unentscheidbarkeit}
%% =============================================================================

\begin{satz}[Boone 1959; Novikov 1955]\label{satz:BN}
Das Wortproblem für endlich präsentierte Gruppen ist unentscheidbar.
\end{satz}

\begin{satz}[Markov, 1958]\label{satz:markov}
Es ist unentscheidbar, ob zwei endliche Präsentationen isomorphe Gruppen definieren.
\end{satz}

\begin{bemerkung}[Implikationen für AC]\label{bem:undec}
Die Unentscheidbarkeit des Isomorphieproblems impliziert nicht direkt die
Unentscheidbarkeit der AC-Vermutung, da diese eine feste Instanz (die triviale Gruppe)
betrifft. Ob die Frage „Ist $P$ AC-äquivalent zur Standardpräsentation?" entscheidbar
ist, ist ebenfalls offen.
\end{bemerkung}

%% =============================================================================
\section{Computersuchergebnisse}
%% =============================================================================

\begin{satz}[Rechnerische Verifikation für kleine Komplexität]\label{satz:comp}
Alle balancierten Präsentationen der trivialen Gruppe mit Gesamtlänge (Summe der
Relationslängen) $\leq 12$ sind AC-äquivalent zur Standardpräsentation.
(Genauer Wert variiert nach Autor; aktuell ca.\ $\leq 14$ verifiziert.)
\end{satz}

\begin{bemerkung}
Havas und Ramsay~\cite{HR03} sowie Bridson~\cite{Bridson15} führten systematische
Computersuchen durch:
\begin{itemize}
  \item Alle Präsentationen niedriger Komplexität reduzieren, aber die
        Reduktionsfolgen können enorm lang sein.
  \item $\mathrm{AK}(2)$ konnte trotz extensiver Suche nicht reduziert werden.
  \item Der Suchraum wächst überexponentiell in der Präsentationslänge.
\end{itemize}
\end{bemerkung}

\begin{proposition}[Längenerhöhende Züge sind notwendig]
Für $\mathrm{AK}(2)$ muss jede AC-Reduktion (falls sie existiert) durch Präsentationen
mit Gesamtrelationslänge $\geq 100$ gehen, was erschöpfende Suche mit heutigen
Methoden undurchführbar macht.
\end{proposition}

%% =============================================================================
\section{Defizit von Gruppen und balancierte Präsentationen}
%% =============================================================================

\begin{definition_de}
Das \emph{Defizit} einer endlichen Gruppe $G$, geschrieben $\mathrm{def}(G)$, ist das
Maximum über alle endlichen Präsentationen $\langle X \mid R \rangle$ von $G$ von
$|X| - |R|$.
\end{definition_de}

\begin{satz}[Defizitschranke]\label{satz:defiz}
Für eine endliche Gruppe $G$: $\mathrm{def}(G) \leq 0$.
Außerdem: $\mathrm{def}(G) = 0$ genau dann, wenn $G$ eine balancierte Präsentation besitzt.
\end{satz}

%% =============================================================================
\section{Verschärfungen und Variationen}
%% =============================================================================

\begin{vermutung}[Stabile AC-Vermutung]\label{verm:stableAC}
Die Andrews--Curtis-Vermutung gilt nach \emph{Stabilisierung}: Für jede balancierte
Präsentation $P$ der trivialen Gruppe gibt es $k \geq 0$, sodass $P * S_k$ AC-reduzierbar ist.
\end{vermutung}

\begin{satz}[AC für Ein-Relationen-Gruppen]\label{satz:einrel}
Die Andrews--Curtis-Vermutung gilt für balancierte Präsentationen der trivialen Gruppe
mit einem Erzeuger und einer Relation (Fall $n=1$): Hier gilt notwendigerweise $r = x^{\pm 1}$,
was bereits Standardform ist.
\end{satz}

%% =============================================================================
\section{Offene Probleme}
%% =============================================================================

\begin{enumerate}[label=(\arabic*)]
  \item \textbf{AC für $n = 2$}: Ist jede balancierte Zwei-Erzeuger-Präsentation
        der trivialen Gruppe AC-reduzierbar? Dies ist der einfachste offene Fall
        und umfasst $\mathrm{AK}(2)$.

  \item \textbf{$\mathrm{AK}(n)$-Präsentationen}: Ist $\mathrm{AK}(n)$ für
        irgendeinen Wert $n \geq 2$ AC-reduzierbar?

  \item \textbf{Topologie}: Besitzt jedes zusammenziehbare $4$-dimensionale
        Handlekörper-Gebäude eine Henkelzerlegung, die aus der Standardzerlegung
        durch eine ableitbare Folge hervorgeht?

  \item \textbf{Algorithmische Entscheidbarkeit}: Ist das Problem „Gegeben $P$,
        ist $P$ AC-äquivalent zur Standardpräsentation?" entscheidbar?

  \item \textbf{Quantitative Schranken}: Falls $P$ AC-reduzierbar ist, wie lang ist die
        kürzeste Reduktionsfolge als Funktion von $|P|$?
\end{enumerate}

%% =============================================================================
\section{Zusammenfassung}
%% =============================================================================

Die Andrews--Curtis-Vermutung steht trotz ihrer elementaren Formulierung als eines
der hartnäckigsten offenen Probleme der kombinatorischen Gruppentheorie und der
niedrigdimensionalen Topologie. Die Gegenbeispielkandidaten $\mathrm{AK}(n)$ und
$\mathrm{MS}(w)$ haben jahrzehntelangen Computersuchen widerstanden, und die Verbindung
zur $4$-Mannigfaltigkeitstopologie legt nahe, dass ein Gegenbeispiel tiefgreifende
Konsequenzen für die Klassifikation glatter $4$-Mannigfaltigkeiten hätte.

\begin{thebibliography}{99}

\bibitem{AC65}
J.~J.\ Andrews, M.~L.\ Curtis,
\textit{Free groups and handlebodies},
Proc.\ Amer.\ Math.\ Soc.\ \textbf{16} (1965), 192--195.

\bibitem{AK79}
S.~Akbulut, R.~Kirby,
\textit{Mazur manifolds},
Michigan Math.\ J.\ \textbf{26} (1979), Nr.\ 3, 259--284.

\bibitem{MS99}
C.~F.\ Miller~III, P.~E.\ Schupp,
\textit{Some presentations of the trivial group},
in \textit{Groups, Languages and Geometry} (R.~Gilman, Hrsg.),
Contemp.\ Math.\ \textbf{250}, AMS (1999), 113--115.

\bibitem{HR03}
G.~Havas, C.~Ramsay,
\textit{Breadth-first search and the Andrews-Curtis conjecture},
Internat.\ J.\ Algebra Comput.\ \textbf{13} (2003), Nr.\ 1, 61--68.

\bibitem{Bridson15}
M.~R.\ Bridson,
\textit{The complexity of balanced presentations and the Andrews-Curtis conjecture},
Preprint, arXiv:1504.04187 (2015).

\bibitem{Smale60}
S.~Smale,
\textit{Generalized Poincaré's conjecture in dimensions greater than four},
Ann.\ of Math.\ (2) \textbf{74} (1961), 391--406.

\bibitem{Tietze08}
H.~Tietze,
\textit{Über die topologischen Invarianten mehrdimensionaler Mannigfaltigkeiten},
Monatsh.\ Math.\ Phys.\ \textbf{19} (1908), 1--118.

\bibitem{LS77}
R.~C.\ Lyndon, P.~E.\ Schupp,
\textit{Kombinatorische Gruppentheorie},
Springer, Berlin, 1977.

\end{thebibliography}

\end{document}
